\section{The \texorpdfstring{$\cxor$}{cxor} architecture}\label{sec:cta}

Let $\F=\{0,1\}$, with addition interpreted as exclusive-or.  Fix a context order $m\ge0$.
The context set is $A=\F^m$, and a table $s:A\to\F$ contains one bit for every context.  Before
the input starts, the table has a known initial value $s_0$ and the context window has a known
prefix $\pi\in\F^m$.  The reference $\W5$ program uses zeros for both.

For an input $x=x_0x_1\dots x_{n-1}$, the address $a_k$ at position $k$ is the length-$m$ suffix
of $\pi x_0\dots x_{k-1}$.  In particular, $a_k$ contains only source symbols \emph{strictly
before} $x_k$.  At a step, write $u=s[a_k]$ for the state read from the addressed cell.  The
general datapath is
\begin{equation}\label{eq:machine-step}
  r_k=x_k+u,
  \qquad
  s[a_k]\leftarrow w(u,x_k),
\end{equation}
where $w$ is a one-bit update rule.  The historical instance uses $w=\W5$, so the second
assignment is simply $s[a_k]\leftarrow x_k$.

\begin{figure}[ht]
  \centering
  \begin{tikzpicture}[
    >={Stealth[length=2.4mm]},
    box/.style={draw, rounded corners=1.5pt, inner xsep=6pt, inner ysep=4pt, font=\small},
    cell/.style={draw, minimum size=6.5mm, inner sep=0pt, font=\footnotesize},
    xorg/.style={draw, circle, minimum size=5.5mm, inner sep=0pt},
    lbl/.style={font=\footnotesize},
    wire/.style={line width=0.5pt},
  ]
  % ---------------- tape
  \node[cell] at (2.35,4.3) {};
  \foreach \p in {3.0,3.65,4.3,4.95} { \node[cell] at (\p,4.3) {}; }
  \node[cell, fill=black!12] (ck) at (5.6,4.3) {};
  \node[cell] at (6.25,4.3) {};
  \node[cell] at (6.9,4.3) {};
  \node[lbl] at (1.6,4.3) {tape $\tape:$};
  \node[lbl, above=2pt of ck] {$\tape(k)$};
  \draw[decorate, decoration={brace, amplitude=4pt, mirror}, wire]
       (2.675,3.9) -- (5.275,3.9)
       node[midway, below=4pt, lbl] {window: the $m$ symbols \emph{before} $k$};

  % ---------------- address
  \node[box] (adr) at (3.975,2.75) {address $\addr{k}$};
  \draw[->,wire] (3.975,3.45) -- (adr.north);

  % ---------------- table
  \node[box, minimum width=2.8cm, minimum height=9mm] (tab) at (2.4,1.3)
       {table $\tab:\addrset\to\F$};
  \draw[->,wire] (adr.south) -- (3.975,2.05) -- (2.4,2.05) -- (tab.north);

  % ---------------- emitter
  \node[xorg] (xr) at (5.6,1.3) {$+$};
  \draw[->,wire] (tab.east) -- node[above, lbl] {$u=\cell{\addr{k}}$} (xr.west);
  \draw[->,wire] (ck.south) -- node[right, lbl, pos=0.66] {$\tape(k)$} (xr.north);
  \node[box] (out) at (7.5,1.3) {$\out(k)$};
  \draw[->,wire] (xr.east) -- (out.west);

  % ---------------- wiring
  \node[box, minimum width=2.4cm, fill=black!5] (w) at (5.6,-0.6) {wiring $\wir(u,\tape(k))$};
  \draw[->,wire] (tab.south) -- (2.4,-0.6) -- node[above,lbl,pos=0.4] {$u$} (w.west);
  \fill (5.6,2.2) circle (1.4pt);
  \draw[->,wire] (5.6,2.2) -- (8.6,2.2) -- (8.6,-0.6) -- (w.east);

  % ---------------- write-back
  \draw[->,wire, line width=0.8pt]
       (w.south) -- (5.6,-1.6) -- node[below, lbl, pos=0.55] {write one bit back to $\cell{\addr{k}}$}
       (0.45,-1.6) -- (0.45,1.3) -- (tab.west);
\end{tikzpicture}

  \caption{One $\cxor$ step.  The preceding $m$ source symbols address a one-bit prediction; the
  exclusive-or emits a residual; and the update rule writes the addressed cell.  In the
  historical $\W5$ implementation the shaded update is $\W5(u,x)=x$.}
  \label{fig:cxor-datapath}
\end{figure}

\begin{algorithm}[ht]
\caption{$\cxor_w$; choose $w=\W5$ for the historical instance}\label{alg:cxor}
\DontPrintSemicolon
$c\leftarrow\pi$; $s\leftarrow s_0$\;
\For{$k\leftarrow0$ \KwTo $n-1$}{
  $a\leftarrow c$; $u\leftarrow s[a]$\;
  $r_k\leftarrow x_k+u$\;
  $s[a]\leftarrow w(u,x_k)$\;
  $c\leftarrow\operatorname{suffix}_m(c x_k)$\;
}
\Return $r$
\end{algorithm}

The ordering matters.  The lookup precedes the window shift, so the prediction is causal.  The
current symbol is used to update the prediction only after its residual has been emitted.

\subsection{Why the reference code is exactly \texorpdfstring{$\W5$}{W5}}

The reference implementation expresses its update conditionally:
\begin{lstlisting}
diff = state != x
out_bits.append(diff)
if diff:
    lookup[index] = x
\end{lstlisting}
For bits, \code{state != x} is $s[a_k]\oplus x_k$, so \code{diff} is precisely the residual
$r_k$.  When \code{diff} is true, the code stores $x_k$.  When it is false, the assignment is
unnecessary because $s[a_k]=x_k$ already.  Consequently the conditional is equivalent to the
unconditional update
\[
  s[a_k]\leftarrow x_k,
\]
which is $\W5$.  Operationally, the program reads the prediction, reports whether it was wrong,
and makes the new prediction equal to the observation.

\subsection{A small trace}

Take $m=2$, a zero prefix and table, and the periodic input $x=\bstr{010101010}$.  The first
occurrences of contexts $\bstr{01}$ and $\bstr{10}$ establish their successors.  Thereafter the
same associations repeat, so their residuals are zero.

\begin{table}[ht]
\centering\small
\begin{tabular}{@{}rccccc@{}}
\toprule
$k$ & context $a_k$ & source $x_k$ & prediction $u$ & residual $r_k$ & new prediction\\
\midrule
0 & \bstr{00} & 0 & 0 & 0 & 0\\
1 & \bstr{00} & 1 & 0 & 1 & 1\\
2 & \bstr{01} & 0 & 0 & 0 & 0\\
3 & \bstr{10} & 1 & 0 & 1 & 1\\
4 & \bstr{01} & 0 & 0 & 0 & 0\\
5 & \bstr{10} & 1 & 1 & 0 & 1\\
6 & \bstr{01} & 0 & 0 & 0 & 0\\
7 & \bstr{10} & 1 & 1 & 0 & 1\\
8 & \bstr{01} & 0 & 0 & 0 & 0\\
\bottomrule
\end{tabular}
\caption{$\W5$ on a two-bit context.  The output is $\bstr{010100000}$.  Each zero is a correct
hard prediction, including initial predictions from the zero table.}
\label{tab:small-w5-trace}
\end{table}

The trace also shows why periodicity is only a special case of the mechanism.  The machine does
not search for a period; it independently remembers that $\bstr{01}$ was followed by $0$ and
$\bstr{10}$ by $1$.

\subsection{Sequential decoder}

The residual loses no information.  A decoder starts with the same prefix, table, and update
rule.  Suppose it has already recovered $x_0,\dots,x_{k-1}$.  Those bits determine $a_k$, and
replaying the same updates determines the encoder's current $u=s[a_k]$.  Hence
\begin{equation}\label{eq:decode-short}
  x_k=r_k+u.
\end{equation}
The decoder then applies $s[a_k]\leftarrow w(u,x_k)$ and advances the window with the recovered
symbol.  This proves the following general fact.

\begin{theorem}[Causal residuals are invertible]\label{thm:causal-inverse}
For every $m$, initial table $s_0$, prefix $\pi$, deterministic update rule $w$, and length $n$,
Algorithm~\ref{alg:cxor} is a bijection $\F^n\to\F^n$.
\end{theorem}

\begin{proof}
Equation~\eqref{eq:decode-short} recovers the input sequentially.  Induction on $k$ shows that
encoder and decoder have identical windows and tables before every step, so the recovery is
unique.  The resulting inverse maps the finite set $\F^n$ onto itself.
\end{proof}

The proof does not invert $w$ and uses no special property of $\W5$.  It inverts the exclusive-or
emitter using a prediction reproducible from the already decoded prefix.  This is the same causal
principle used by predictive lossless codes with much richer predictors.

If the window were shifted before lookup, so that $x_k$ participated in its own address, the
decoder would in general need $x_k$ to discover the prediction required to recover $x_k$.
Collisions then occur for ordinary nonzero tables.  The lagging, source-driven address in
Algorithm~\ref{alg:cxor} is therefore part of the transform, not a cosmetic indexing convention.
